\chapter{Applications, Advantages and Limitations}
\label{ch:adv}
\label{ch:apps}

\section{Applications}

The design choices of \NS---a million-token context, high throughput on long
generations, and post-training focused on multi-step tool use---point to a
particular class of applications. The Nemotron~3 white paper positions Super
for collaborative agents and high-volume workloads such as IT ticket
automation~\cite{nvidia2025nemotron3}.

\begin{itemize}
  \item \textbf{Coding agents.} Reading a repository, editing files and
        running tests over many steps produces long, prefill-heavy contexts.
        \NS\ was trained on real GitHub issues and scores 60.5 on SWE-Bench.
  \item \textbf{Terminal and tool-using agents.} Command-line automation,
        database queries, web search and customer-service tool calls
        (TauBench) benefit from high throughput and a cost per token that does
        not climb as the session grows.
  \item \textbf{Long-document analysis.} Entire contracts, research papers or
        codebases fit in a one-million-token context, and the RULER score of
        91.6 shows the model can actually use it.
  \item \textbf{Reasoning with controllable effort.} Three reasoning
        modes---off, regular and low-effort---combined with a reasoning budget
        let an application spend long traces on hard problems and answer simple
        requests quickly, trading accuracy against cost per request.
  \item \textbf{On-premises deployment.} Open weights and FP8/NVFP4
        checkpoints let organisations in finance, healthcare or government run
        and fine-tune the model on their own hardware.
\end{itemize}

\section{Advantages}

\begin{enumerate}
  \item \textbf{Long context at low cost:} only 8 of 48 sequence-mixing layers
        need a KV cache; the rest keep a constant-size state.
  \item \textbf{High throughput:} up to 2.2$\times$ and 7.5$\times$ faster
        than GPT-OSS-120B and Qwen3.5-122B at comparable accuracy.
  \item \textbf{Large capacity, small active size:} about 12B of 120B
        parameters are used per token.
  \item \textbf{Hardware-aware MoE:} LatentMoE cuts expert memory and routing
        traffic fourfold and spends the saving on more experts.
  \item \textbf{Built-in speculative decoding:} MTP heads accept 3.45 tokens
        per step on average, with no separate draft model.
  \item \textbf{Openness:} weights, much of the data and the recipe are
        public.
\end{enumerate}

\section{Limitations}

\begin{enumerate}
  \item \textbf{Memory footprint:} all 120B parameters must be held in memory.
  \item \textbf{Serving complexity:} expert routing across GPUs, and support
        for Mamba, attention, MoE and MTP together, complicate deployment.
  \item \textbf{Sensitive recurrent state:} quantizing the Mamba state causes
        drift; the fix relies on recent NVIDIA hardware.
  \item \textbf{Hardware-dependent speed:} the 2.2$\times$ figure needs native
        FP4 support.
  \item \textbf{Accuracy gap:} it trails Qwen3.5-122B on several reasoning and
        agentic benchmarks.
  \item \textbf{Vendor-reported results:} all figures come from the developers;
        independent evaluation is still needed.
\end{enumerate}

\section{Comparison at a Glance}

Table~\ref{tab:compare} places \NS\ beside a pure Transformer MoE model of the
same class.

\begin{table}[!htbp]
\centering
\small
\caption{Hybrid Mamba--Transformer MoE compared with a pure Transformer MoE}
\label{tab:compare}
\begin{tabularx}{\linewidth}{@{}P{3.6cm}LL@{}}
\toprule
\textbf{Aspect} & \textbf{Pure Transformer MoE} & \textbf{\NS\ (hybrid)} \\
\midrule
Sequence mixing & Attention in every block & Mamba-2 in 40 blocks, attention in 8 \\
Generation memory & KV cache in every block & KV cache in 8 layers only \\
Feed-forward part & Standard MoE at full width & LatentMoE in a 1024-dim latent space \\
Faster decoding & External draft model, if any & Built-in MTP heads \\
Deployment risk & Well understood & Recurrent state needs careful quantization \\
Tooling maturity & Very mature & Newer, fewer optimised runtimes \\
\bottomrule
\end{tabularx}
\end{table}
